\section{第四讲：MPO 代数：从代数 Bethe ansatz 到 G-单射 PEPS}
现在让我们将注意力转向可积模型（integrable models），这是本讲习班的核心主题之一。事实证明，张量网络提供了一种非常简明的方式，可以用矩阵乘积态与算符来重新表述\emph{代数 Bethe ansatz}（algebraic Bethe ansatz, ABA）。特别是，我们将看到构成其可积性核心的 \emph{Yang-Baxter 方程}（Yang-Baxter equation），可以被理解为基本定理（\cref{sec:FundamentalTheorem}）的一个推论。尽管张量网络技术并未向我们揭示关于 Bethe ansatz 的任何根本全新的物理，但本讲的见解将对由\emph{融合范畴}（fusion categories）所描述的\emph{广义}或\emph{范畴对称性}（categorical symmetries）的 MPO 表示论发挥关键作用，而这正是最后一讲的主题。

在讨论完代数 Bethe ansatz 之后，我们将在\cref{sec:PEPS}中重点讨论 MPS 的二维推广，即\emph{投影纠缠对态}（projected entangled-pair states, PEPS）。我们将猜想基本定理的一种推广形式，并讨论 MPO 对称性在二维 SPT 与拓扑相及其边缘模的张量网络表示中所起的作用。
\subsection{磁振子态}\label{sec:magnon}
张量网络可以用作 Hamiltonian 解析解的 ansatz，正如针对 AKLT 模型\cref{eq:AKLTState}所展示的那样。然而，本讲所关注的范例是一类著名的精确可解模型，称为\emph{可积系统}（integrable systems）。粗略地说，可积量子系统是具有广延数目的守恒荷（conserved charges）——即与 Hamiltonian 对易的局域算符——的系统，其精确可解性正归功于此。可积系统可以用称为 \emph{Bethe ansatz} 的 ansatz 波函数来求解（参见文献~\cite{Gaudin2014}）。该方法有两种表述形式，分别称为\emph{坐标 Bethe ansatz}（coordinate Bethe ansatz）与\emph{代数 Bethe ansatz}（algebraic Bethe ansatz），它们产生同一组解。这两种表述形式都可以使用张量网络写出。

为了具体起见，我们考虑自旋 $1/2$ XXZ Hamiltonian
\begin{equation}
    H_{XXZ}=\sum_{\msf i=}^L \sigma^X_{\msf i}\sigma^X_{\msf i+1}+\sigma^Y_\msf i\sigma^Y_{\msf i+1}+\Delta \sigma^Z_\msf i\sigma^Z_{\msf i+1}
\end{equation}
在周期性边界条件（即 $L+1\equiv 1$）下。求解 Bethe 方程时用作基底的\emph{磁振子态}（magnon states）可以很容易地写为 MPS。例如，\emph{周期性}边界条件下的 1-磁振子态可以写为\emph{开}边界条件 MPS（参见\cref{eq:OpenMPS}），其左边界矢量为 $\vec v_L^\dagger = \bra{0}$，右边界矢量为 $\vec v_R = \ket{1}$，MPS 张量为：
\begin{equation}
\begin{split}
A^0(k) = 
\begin{pmatrix}
    e^{ik} & 0 \\
    0 & 1
\end{pmatrix}, \quad
A^1(k) = \begin{pmatrix}
    0 & e^{ik} \\
    0 & 0
\end{pmatrix} \\
\implies \ket{\Psi_1[A]} = \sum_{\msf n=1}^L e^{ik\msf n} \ket{0_0 ...0_{\msf n-1} 1_\msf n 0_{\msf n+1}...0_{L}}
\end{split}.
\label{AAAA}
\end{equation}
注意，这正是在\cref{sec:Excitations}中在 MPS 流形上将激发态构造为具有动量的切矢量的方式。一般 n-磁振子态可以在文献~\cite{Alcaraz2004,Murg2012,Haegeman2016}中找到，并涉及一个键维随系统中粒子数呈指数级增长的 MPS。
\subsection{作为 MPS 基本定理的 Yang-Baxter 方程}\label{sec:YangBaxterFromMPO}
然而，还有一种完全不同的方法，即借助代数 Bethe ansatz（ABA）导出精确本征态。在用张量网络表述这两种方法之前，并不清楚这两种方法彼此之间有何关联；现在人们已经清楚，正如由 MPS 基本定理所必然要求的那样，这两种方法通过一个非平庸的 gauge 变换相关联~\cite{Katsura2010}。在代数 Bethe ansatz 中，Hamiltonian 的所有本征态都可以通过依赖于连续\emph{谱参数}（spectral parameter）$\mu_i$ 的 MPO 相乘来构造：
\begin{equation}\label{eq:BetheAnsatzMPS}
    \inputtikz{YB_ABA}.
\end{equation}
此处，所有 MPO 张量均具有如下形式：
\begin{equation}
    A_{\alpha\beta}^{ij}(\mu) =\mu\delta_{\alpha\beta}\delta_{ij}+\delta_{\alpha i}\delta_{\beta j}= \mu \inputtikz{YB_A_1} + \inputtikz{YB_A_2}.
    \label{AA}
\end{equation}
容易检验，每个 MPO 都会向系统中增加一个磁振子，因为左边界为 $\bra{0}$，右边界为 $\ket{1}$。系数 $\mu_i$ 是磁振子动量的函数，可以通过求解所谓的 \emph{Bethe 方程}（Bethe equations）来确定。ABA 的一个关键要素是所有这些 MPO 均相互对易。

其背后的逻辑如下。我们从一个二维统计力学模型的转移矩阵出发，该转移矩阵可以写为单射 MPO $T(\mu)$ 的形式，由与式~\ref{AA}~中相同的张量 $A(\mu)$ 确定，但此时具有周期性边界条件。事实证明，该系统可积当且仅当所有这些 MPO 与其谱参数无关地相互对易：
\begin{equation}
        \forall \lambda, \mu: [T(\lambda), T(\mu)] = 0.
\end{equation}
在此情形下，我们可以调用 MPS 基本定理\cref{eq:FundTh}来证明存在一个将这些乘积相互转化的 gauge 变换 $R$~\cite{Haegeman2016}。事实上，注意我们可以将 MPO 的两条腿合并为一条腿，从而可以将 MPS 基本定理用于秩为 4 的 MPO。具体而言，存在一个满足如下方程的交织算符（intertwiner）张量 $R(\lambda, \mu)$：
\begin{equation}\label{eq:YB_intertwiner}
    \inputtikz{YB_intertwiner_1} \,=\, \inputtikz{YB_intertwiner_2}\,\,.
\end{equation}
该 $R$-张量是由对易 MPO 给出的全局对称性的局域实现。

现在让我们取 $R$-张量的三个副本。由于 MPO 代数的结合律，这些 $R$-张量可以通过两种不同方式融合在一起，从而导出如下方程：
\begin{equation}
    \inputtikz{YB_1} \,=\, \inputtikz{YB_2}\,\,.
\end{equation}
这正是 Yang-Baxter 方程——因此 MPS 基本定理以非常简明的方式重现了这一结果。给定任意可积转移矩阵，相应的交织算符因此必须满足 Yang-Baxter 方程。可以检验，\cref{AA} 中的张量正是这些解之一。因此，寻找可积模型的一种构造性方法是将思路反转过来：首先寻找 Yang-Baxter 方程的一个解（这通常极其困难！），从而得出张量 $R(\lambda,\mu)$ 的显式表示。下一步，将转移矩阵 $T$ 作为\cref{eq:YB_intertwiner}的解求解出来。令人惊叹的是，通过将 MPO 张量识别为 $A_{\alpha \beta}^{ij}(\mu) := R_{j \beta}^{i \alpha}(\mu, 0)$，该方程等价于 Yang-Baxter 方程本身。因此，Yang-Baxter 代数的这一所谓\emph{基本表示}（fundamental representation）允许构造单参数族对易转移矩阵。注意，这并不是求解\cref{eq:YB_intertwiner}的唯一 MPO 族。人们可以构造同一代数的不同表示——在后文中，我们将遇到类似的情形，届时我们将构造同一对称系统的不同表示。作为最后一步，现在可以用相同的张量 $A$ 但采用不同的边界条件，来构造\cref{eq:BetheAnsatzMPS}中的磁振子 MPO。

因此，张量网络表述的核心原理与优势如下：\emph{全局}守恒律和对易关系可以写成涉及\emph{局域}自由度（即 MPS 和 MPO 张量）的方程，而这些方程通常更容易求解。

在给出 MPO 表述之后，还可以恢复可积 Heisenberg Hamiltonian。所有权重彼此相等的 6-顶点模型（相应的可积统计力学模型）的转移矩阵恰好具有\cref{AA}的形式。根据构造，具有这些张量 $A(\mu)$ 的所有 MPO $T(\mu)$ 都将相互对易，因此该 MPO 的导数是唯一确定的。Heisenberg Hamiltonian 确实是该转移矩阵的对数导数；我们将证明留作一个（非常有趣的）练习（详尽讨论参见文献~\cite{Murg2012}）：
\begin{equation}
    H = T^{-1}(0) \frac{{\rm d}T(\mu)}{{\rm d}\mu}\vert_{\mu=0} = \sum_{\msf i}\quad\inputtikz{YB_Ham}\,.
\end{equation}

总之，两种 Bethe ansatz 表述（坐标和代数）在张量网络中都有自然的表示。张量网络与 Bethe ansatz 显然共享着相同的基因，但张量网络已被证明具有更广泛的适用性——它们同样适用于不可积系统和更高维度。
\subsection{投影纠缠对态}\label{sec:PEPS}
矩阵乘积态向二维的自然推广是所谓的\emph{投影纠缠对态}（projected entangled-pair states, PEPS）~\cite{Verstraete2004b}。与一维情形类似，平移不变的 PEPS 由单个张量构造而成，在正方晶格的情形下，该张量具有一个物理指标和四个虚指标：
\begin{equation}
    A = \sum_{\substack{i,\alpha,\beta\\\gamma,\delta}} A^i_{\alpha\beta\gamma\delta} \ket{i}\otimes\bra{\alpha}\otimes\bra{\beta}\otimes\bra{\gamma}\otimes\bra{\delta} \in \mbb C^{d\chi^4}.
\end{equation}
该张量如图所示：
\begin{equation}
    \inputtikz{PEPSTensor} = \ A^i_{\alpha\beta\gamma\delta}, \qquad \substack{i=1,2,\dots,d\\ \alpha,\beta,\gamma,\delta=1,2,\dots,\chi},
\end{equation}
从而整个 PEPS 波函数可以表示为
\begin{equation}
    \ket{\Psi[A]} = \inputtikz{PEPS}\,\, .
\end{equation}
注意，PEPS 自然满足面积律（area law），因为一段连续自旋区域的纠缠熵随该区域边界上的键数（即其周长）线性增长。因此，PEPS 预期为局域有能隙 Hamiltonian 的低能态构成一个良好的变分族，而长期以来人们一直猜想这些低能态满足面积律。反之，与 MPS 一样，可以证明每个 PEPS 都是某个局域无阻挫（frustration free）有能隙 parent Hamiltonian 的精确基态。如何将这种构造推广到开边界条件、任意晶格或更高维度也是显而易见的。类似地，人们可以将 MPO 推广到二维情形。此类算符被称为\emph{投影纠缠对算符}（projected entangled-pair operators, PEPO）。

人们付出了相当大的努力来将 DMRG 和 VUMPS 算法推广到 PEPS 的情形。虽然原则上将流形图像推广到 PEPS 是直截了当的，但其关于键维的实际计算复杂度增长远差于一维情形。尽管如此，这些 PEPS 算法仍在二维强关联晶格系统（例如 Hubbard 模型）中提供了最前沿的变分结果，提出更好的算法是一个非常活跃的研究领域。关于最前沿的结果，我们推荐参阅文献~\cite{Verstraete2008,Vanderstraeten2016,Corboz2016,Liao2019,Vanderstraeten2022}。在这些讲义中，我们将侧重于 PEPS 的理论。

与矩阵乘积态的情形相反，目前尚不存在完全通用的基本定理，能够为两个不同的 PEPS 张量 $A$ 和 $B$ 在所有系统尺寸下生成相同的态提供充分必要条件。然而很明显，使其成立的一个充分条件是存在一个矩阵乘积算符，按照如下方式\emph{交织}（intertwines）PEPS 张量：
\begin{equation}
    \inputtikz{PEPSPullingThrough_1} \,\,=\,\,
    \inputtikz{PEPSPullingThrough_2}\,\,,
\end{equation}
以及所有其他穿透（pulling through）MPO 的方式，例如：
\begin{equation}
    \inputtikz{PEPSPullingThrough_3} \,\,=\,\,
    \inputtikz{PEPSPullingThrough_4}\,\,.
\end{equation}
在缺乏严格证明的情况下，我们称其为 PEPS 的基本\emph{猜想}（conjecture）。假设周期性边界条件，对称性作用的实现方式是：从 PEPS $\ket{\Psi[A]}$ 出发，成核一个由交织 MPO 构成的小泡，将其穿透整个晶格并与自身湮灭，最终得到 PEPS $\ket{\Psi[B]}$。正如我们将在本讲义后续部分所看到的，大类 PEPS 恰好是按照这一方案来实现其对称性的。在某些情况下，该 MPO 只是一个直积（键维为 $1$），从而虚对称性仅作为直积作用在虚腿上。

在接下来的各节中，我们将讨论若干表示拓扑物态的特别有趣的 PEPS 波函数。拓扑波函数可以使用张量网络形式轻松构造，这显然是它们受到如此广泛研究的主要原因之一。此外，这些张量网络表示将长程纠缠与作用在虚自由度上的（MPO）对称性等同起来。这也是该基本猜想如此重要的原因：拓扑序由非平庸的 MPO 对称性所表征。
\subsection{二维 SPT：案例研究}\label{sec:SPT2d}
具有非平庸 MPO 对称性的 PEPS 的典型范例是 Chen、Liu 和 Wen 在文献~\cite{Chen2011}中引入的所谓 \emph{CZX 模型}（其张量网络表示参见文献~\cite{Buerschaper2013,Williamson2016,Molnar2017}）。他们的模型是由 $\mbb Z_2$ 对称性保护的非平庸 SPT 的一个例子，该对称性在 PEPS 的物理层面上以单格点方式 $U_g^{\otimes N}$ 起作用。在 PEPS 张量的虚指标层面上，该对称性通过键维为 2 的 MPO 实现为：
\begin{equation}\label{eq:PEPS_CZX}
    \inputtikz{PEPS_CZX_1} = 
    \inputtikz{PEPS_CZX_2}\,\, ,
\end{equation}
其中 $g\in\{0,1\}$。由 $g$ 索引的 MPO 构成了 $\mbb Z_2$ 的非局域表示，其图形表示为
\begin{equation}\label{eq:GroupMPO}
    \inputtikz{MPOGroupStacked} =\,
    \inputtikz{MPOGroupgh}\,\, .
\end{equation}
乘法\cref{eq:GroupMPO}可以通过满足如下条件的交织张量在局域上实现\footnote{$g=0$ 对应的 MPO 不是单射的，但不能分解为单射块。这可以通过底层的非半单代数结构来解释~\cite{Liu2025,FloridoLlinas2025}。因此，无法直接通过调用基本定理来找到交织算符。}~\cite{Chen2011,Williamson2016,Liu2025}：
\begin{equation}\label{eq:GroupIntertwiner}
    \inputtikz{GroupMPOIntertwiner_1} = \inputtikz{GroupMPOIntertwiner_2}\,\, .
\end{equation}
回想在\cref{sec:SPT1d}中我们曾论证，在一维 SPT 相的情形下，存在一个取值于第二上同调群 $H^2(G,\rU(1))$ 的拓扑指数。该指数完全由构成对称性射影表示的虚矩阵的乘法规则所捕获。现在很自然地会想知道二维 SPT 相对应的拓扑指数是什么，以及它是如何编码在 MPO 代数中的。事实证明，如果考虑表示 $\mbb Z_2$ 对称性的三个 MPO 的乘法，那么上面定义的融合张量可以通过两种不同方式进行收缩。这两种方式必须相互兼容，因此必然存在一个依赖于三个群元素的相位 $\omega$，满足
\begin{equation}\label{eq:GroupMPOFusion}
    \inputtikz{GroupMPOPentagon_1}
    = \omega(g,h,k) \,\,
    \inputtikz{GroupMPOPentagon_2}\,\, .
\end{equation}
矩阵乘法的结合律表明 $\omega$ 必须满足
\begin{equation}
    \omega(g,h,k)\omega(g,hk,l)\omega(h,k,l) = \omega(gh,k,l)\omega(g,h,kl).
\end{equation}
该条件被称为 \emph{3-上循环方程}（3-cocycle equation）（亦参见\cref{sec:GroupCohomology}）。注意，我们可以给融合张量自由地乘以一个相位 $\varphi(g,h)$，在此重新定义下，3-上循环按照如下规则变化：
\begin{equation}\label{eq:3cocycle}
    \omega(g,h,k) \mapsto \omega(g,h,k) \frac{\varphi(g,h)\varphi(gh,k)}{\varphi(g,hk)\varphi(h,k)}.
\end{equation}
模去此类重新定义后，我们可以得出结论：指数 $\omega$ 取值于第三上同调群 $H^3(G,\rU(1))$。该群可以通过\cref{sec:GroupCohomology}中给出的方法进行计算，并且特别地它也是有限群。对于上述考虑的 $G=\mbb Z_2$ 情形，有 $H^3(\mbb Z_2,\rU(1))\simeq \mbb Z_2$，因此恰好存在一个由 $\mbb Z_2$ 对称性保护的非平庸 SPT 相。相应非平庸类的代表元 3-上循环为
\begin{equation}
    \omega(g,h,k) = (-1)^{ghk},
\end{equation}
这可以直接从 CZX 模型的融合张量的\cref{eq:GroupIntertwiner}中推导出来。正如在量子自旋链的情形中一样，具有不同虚对称性/3-上循环的两个 PEPS 波函数无法在不闭合能隙的情况下通过保持对称性的绝热路径相互转化。
\subsection{边缘 Hamiltonian 与异常}
当 PEPS 定义在开面片 $\mc A$ 上时，出现了一个值得注意的特征。在这种情况下，未收缩的键定义了一个维度为 $\chi^{\lvert\partial \mc A\rvert}$ 的希尔伯特空间，其中 $\lvert\partial\! \mc A\rvert$ 是与边界相交的连键数目。图形表示为：
\begin{equation}
    \inputtikz{PEPS_open}\,\,.
\end{equation}
这些边缘自由度为相应 parent Hamiltonian 的边缘模定义了一个张量积结构，这与我们在\cref{sec:SPT1d}中识别一维情形下 SPT 边缘模的方式完全一致。当 PEPS parent Hamiltonian 随后受到微扰时，这会在这些边界自由度上诱导出一个\emph{有效} Hamiltonian，并有可能驱动它们经历相变~\cite{Yang2014}。有趣的是，在 SPT 的 PEPS 表示情形下，这些有效边界 Hamiltonian 继承了体态的非局域 MPO 对称性。有关若干示例及其数值研究，请参阅文献~\cite{Williamson2016,Roose2018}。

给出一个群的 MPO 表示，反过来提出了对称（边缘）Hamiltonian 的相图会是什么样的问题。正如我们现在将要展示的，这个问题的答案强烈取决于与 MPO 代数相关联的 3-上循环是否平庸。正如文献~\cite{Chen2011,Cirac2020}中所证明的，一个非常优美且直接的张量网络论证表明，只要 $\omega$ 在上同调上是非平庸的，与之对易的 Hamiltonian 就必须要么自发破缺对称性，要么是无能隙的。在这种情况下，称上循环 $\omega$ 代表了对称性的\emph{异常}（anomaly）。这一观察与著名的 Lieb-Schultz-Mattis 定理~\cite{Lieb1961,Oshikawa2000,Hastings2003}高度平行。论证过程如下。

在该论证中，我们假设任意群 $G$ 的一个\emph{单射} MPO 表示。由于单射性，我们现在可以调用基本定理来证明存在局域实现群乘法的融合张量，如\cref{eq:GroupIntertwiner}所示。现在假设按\cref{eq:GroupMPOFusion}定义的 3-上循环是非平庸的，并假设存在一个由单射 MPS 表示的唯一的、平庸有能隙的对称基态。在图形上，这归结为：
\begin{equation}
    \inputtikz{GroupMPOMPS} \,\, = \,\,
    \inputtikz{GroupMPS} \,\, .
\end{equation}
随后应用 MPS 基本定理表明，必然存在一个满足如下条件的三价\emph{作用张量}（灰色）：
\begin{equation}
    \inputtikz{GroupMPSIntertwiner_1} \,\, = \,\,
    \inputtikz{GroupMPSIntertwiner_2} \,\, .
\end{equation}
与上一节类似，如果我们在态上作用两个对称 MPO，则融合 MPO 张量随后作用在 MPS 上的结果，在相差一个相位的意义下等于相继作用这些 MPO。具体而言：
\begin{equation}
    \inputtikz{GroupMPSIntertwinerRecoupling_1} \,\, = \,\, \lambda(g,h)
    \inputtikz{GroupMPSIntertwinerRecoupling_2} \,\, .
\end{equation}
如果我们现在考虑三个此类 MPO 在我们的 MPS 上的作用，并考虑在局域重新耦合它们的两种不同方式，我们便可以推导出如下条件：
\begin{equation}
    \frac{\lambda(g,h)\lambda(gh,k)}{\lambda(h,k)\lambda(g,hk)}=\omega(g,h,k),
\end{equation}
但左边可以被识别为一个上边缘，这反过来意味着 3-上循环在上同调上是平庸的。由于这与 $\omega$ 是非平庸的假设相矛盾，我们因此必须得出结论：不存在对称单射 MPS，并且每个具有 MPO 对称性的一维 Hamiltonian 都必须是对称性破缺的或无能隙的。注意，无论该 Hamiltonian 编码的是二维 SPT 相的边缘模还是真实的量子自旋链，该论证均成立。
\subsection{拓扑序：量子双重模型与管代数}\label{sec:QuantumDoubles}
引人注目的是，除 SPT 相之外，在高于一维的空间维度中还存在不受任何对称性保护的鲁棒量子物态。称此类相展现出\emph{内禀拓扑序}（intrinsic topological order）~\cite{Wen1989,Witten1988,Zeng2015,Wen2016}，其特征在于仅依赖于底层空间流形拓扑（球面、环面等）的基态简并度、具有\emph{任意子}编织统计的有能隙激发，以及根据其任意子内容决定的对纠缠熵面积律的普遍负修正~\cite{Kitaev1997,Kitaev2005a,Levin2006}。值得注意的是，与不存在对称性紧密相关的是，这些相不能由局域序参量来表征，从而超出了对称性破缺的传统 Landau 范式的范畴~\cite{Landau1937,Wegner1971,Bravyi2006}。尽管如此，存在基态简并度的事实清楚地表明某种对称性被破缺了。事实证明，这种对称性纯粹作用在虚自由度上。

回想短程纠缠的 SPT 序的稳定性源于与全局对称性下纠缠自由度的变换性质相关联的拓扑指数的存在（\cref{sec:SPT1d} 和 \ref{sec:SPT2d}）。另一方面，内禀拓扑序的鲁棒性源于其长程纠缠~\cite{Chen2010}。在某种意义上，表征拓扑序因此归结为对长程纠缠的等价类进行分类。类似地，这些拓扑相之间的相边界可以从纠缠性质的根本性改变来理解~\cite{Haegeman2015,Marien2016}。正如在 PEPS 中那样能够直接处理这些纠缠自由度，这启示我们可以进一步将该问题约化为对产生\emph{全局}拓扑态的\emph{局域} PEPS 张量进行分类：内禀拓扑序系统的分类等价于构造具有 MPO 对称性的张量。

事实上，用 PEPS 的语言来说，拓扑序的存在可以追溯到与基态 PEPS 张量对易的\emph{虚} MPO 对称性的存在：
\begin{equation}\label{eq:PEPSPullingThrough}
    \inputtikz{PEPS_TopoPulling_1} \,=\, \inputtikz{PEPS_TopoPulling_2}\,\,,
\end{equation}
注意，与 SPT 情形相反，在物理自由度上没有任何作用。因此，这些 MPO 在虚指标层面上是完全可形变的——即\emph{拓扑的}。这也意味着底层的对称性结构可以比群所描述的结构丰富得多。正如我们将在第~\ref{sec:GenSym}~讲中所看到的，对称代数由一个\emph{融合范畴}（fusion category）编码。由于这些 MPO 对物理自由度不敏感，此类态对物理层面上的（微小）局域微扰具有鲁棒性~\cite{Marien2016}。然而，\emph{任意子凝聚}（anyon condensation）提供了一种在足够强的微扰下驱动系统走向平庸相的机制~\cite{Bais2008,Burnell2017}。

人们猜想，非手征拓扑序完全由其任意子内容以及融合规则与编织统计所表征。在数学上，这些数据被封装在一个称为 \emph{Drinfeld 中心} $\mc Z(\mc D)$ 的\emph{模张量范畴}（modular tensor category）中~\cite{Etingof2015,Kong2017,Bultinck2015,Kong2022}。该论断的完整范畴以及一般拓扑序的微观实现将是\cref{sec:SN}的主题。在本节中，我们将重点关注一类称为（无扭）\emph{量子双重模型}（quantum doubles）~\cite{Dijkgraaf1989,Roche1990,Kitaev1997}的拓扑序的构造，作为对拓扑序的 PEPS 表示及其若干性质的更为温和的引介。

\bigskip\noindent Kitaev 在文献~\cite{Kitaev1997}中引入的量子双重模型是由有限群 $G$ 构造的一族模型，方法是在二维晶格的边上赋予 $\lvert G \rvert$ 维 qudit，并服从上述文献中详述的对易投影算符 Hamiltonian。它们是无扭 Dijkgraaf-Witten 拓扑场论的晶格 Hamiltonian 实现~\cite{Dijkgraaf1989}。对于群 $\mbb Z_2$，该模型约化为\emph{复环面码}（toric code）。这些量子双重模型的基态允许非常简单的重整化群不动点 PEPS 表示，称为 \emph{$G$-单射}（G-injective），或者更精确地称为 \emph{$G$-等距}（G-isometric）PEPS。$G$-等距的概念本质上抓住了仅在群 $G$ 作用下不变的虚自由度子空间上可逆的 PEPS 张量，是 MPO 单射性的一种具体体现。

在局域物理变换的意义下，所有 $G$-等距 PEPS 都可以约化为我们现在将要描述的普适形式。在本节中，我们假设二维环面上的正方晶格。我们用 $L$ 表示 $G$ 的\emph{左 regular 表示}（left regular representation），其作用为 $L_g\ket{h} = \ket{gh}$。然后定义投影算符 $1/\lvert G \rvert\sum_g L_g\otimes L_g\otimes L_g^\dagger\otimes L_g^\dagger$，它构成了基态 PEPS 张量，图形表示为：
\begin{equation}\label{eq:PEPSGTensor}
    \frac{1}{\lvert G \rvert}\sum_{g\in G}\,\,\, \inputtikz{PEPS_G_1}\,\,.
\end{equation}
在此图中，灰色方框代表 $L_g$，横线表示逆群元素（回想 $L_g^\dagger = L_{\bar g}$）。注意，最内侧的四条腿共同构成了物理层面。由于式~\ref{eq:PEPSGTensor}~中的求和遍历所有群元素，该 PEPS 张量对于如下由群元素标记的 MPO 满足穿透条件\cref{eq:PEPSPullingThrough}：
\begin{equation}\label{eq:MPOGroupRegular}
    \inputtikz{MPOGroupg}:= \bigotimes_\msf i L_g,
\end{equation}
其中虚线用于强调键维为 1，即每个 MPO 均表现为直积。

PEPS 张量\cref{eq:PEPSGTensor}仅给出了该模型的一个基态。正如文献~\cite{Schuch2010}中所展示的，所有其他基态都可以通过用环绕环面任一或两个不可缩回圈的算符作用在该态上而幺正地获得。这些算符的广延性质反映了基态的局域不可区分性。等价地，可以通过在虚指标层面上用 MPO 环绕环面的两个圈并在其交叉点处用一个需要求解的张量 $\mc P_Z$ 将它们连接起来，从而获得所有基态。在此过程中，我们要求基态均具有平移不变性，并在热力学极限下保持可归一化和正交~\cite{Zhang2011}。利用 MPO \cref{eq:MPOGroupRegular} 的融合张量，我们可以将这样的态展开如下：\footnote{注意，此时的融合张量仅仅强制群乘法，即 $\sum_{g,h,k}\ketbra{k}{g,h}=\sum_{g,h}\ketbra{gh}{g,h}$。}
\begin{equation}
    \inputtikz{PEPS_G_idempotent_1} \,\, = \sum_{g,h} \mc (\mc P_Z)_{g,h}\,\,\inputtikz{PEPS_G_idempotent_2}\,\,.
\end{equation}
这些态的平移不变性由穿透条件显式保证。正如文献~\cite{Sahinoglu2014,Bultinck2015,Williamson2017}中所论证的，MPO 单射 PEPS 中的一组基态基底与 \emph{Ocneanu 管代数}（tube algebra）~\cite{Ocneanu1993,Ocneanu2001}的\emph{极小中心幂等元}（minimal central idempotents）——或者等价地，不可约表示——一一对应。事实上，在实践中 Drinfeld 中心 $\mc Z(\mc D)$ 通常就是通过管代数计算的。对于量子双重模型的特定情形，相关的管代数由如下形式的闭合 MPO 张成：
\begin{equation}
    T_x^g := \inputtikz{PEPS_G_Tube}\,\, ,
\end{equation}
可将其视为从外侧指标作用到内侧指标。它们的乘法定义为将它们相互套叠。在代数上这归结为：
\begin{equation}\label{eq:GTubesMult}
    T^g_x\circ T_y^h = \delta_{x,hyh^{-1}} T^{gh}_x.
\end{equation}
该乘法定义了一个通常称为群 $G$ 的\emph{（无扭）量子双重模型}（quantum double）$\mc D(G)$ 的代数~\cite{Roche1990}。值得注意的是，管代数赋予了一个 \emph{star} 结构，即交换乘积顺序的反线性对合。对于当前情形：$\big( T_x^g\big)^\star := T^{g^{-1}}_{gxg^{-1}}$。
因此，管代数被赋予了 \emph{$C^\star$-代数}的结构。根据著名的 Artin-Wedderburn 定理，每个 $C^\star$-代数都可以分解为单矩阵代数。前面提到的极小中心幂等元便充当投射到每个单块上的投影算符。

让我们显式地计算它们。为此，对于固定的 $x\in G$，首先关注所有满足 $g\in Z_x=\{g\in G \mid gx=xg\}\subseteq G$ 的管 $T_x^g$ 是非常有益的。这些管构成了稳定子群 $Z_x$ 的一个表示，因为
\begin{equation}
    T^{g_1}_x\circ T^{g_2}_x = T^{g_1g_2}_x, \quad \forall g_1,g_2\in Z_x,
\end{equation}
正如直接由\cref{eq:GTubesMult}所给出的那样。令 $\mu_x$ 表示 $Z_x$ 的一个维度为 $d_{\mu_x}$ 且相应特征标为 $\chi^{\mu_x}$ 的不可约表示。算符
\begin{equation}
    \mc P_{(x,\mu)} := \frac{d_{\mu_x}}{\lvert G \rvert} \sum_{k\in Z_x} \chi^{\mu_x}(k) T_x^k,
\end{equation}
构成了 Hermitian 投影算符，即\footnote{利用特征标的 Schur 正交性 $\sum_k \chi^{\mu_x}(k)\chi^{\nu_x}(k^{-1}m) = \lvert Z_x \rvert/d_\mu \chi^\mu(m)$。}
\begin{equation}
    \mc P_{(x,\mu_x)} \mc P_{(y,\nu_y)} = \delta_{x,y}\delta_{\mu_x,\nu_x} \mc P_{(x,\mu_x)}.
\end{equation}
这些 $\mc P_{(x,\mu)}$ 是\emph{单}幂等元（simple idempotents）。中心的幂等元随后通过对共轭类 $C_x = \{gxg^{-1} \mid g\in G\}$ 中的所有 $x$ 额外求和而获得：
\begin{equation}
    \mc P_{(C_x,\mu_x)} \equiv \sum_{y\in C_x} \mc P_{(y,\mu_x)}.
\end{equation}
这重现了文献~\cite{Kitaev2005a}的结果，即量子双重模型的环面基态由共轭类 $C_x$ 以及 $C_x$ 的代表元 $x$ 的稳定子子群 $Z_x$ 的不可约表示所标记。

重要的是，任意子激发也由管代数的极小中心幂等元所标记~\cite{Lan2013,Hu2015}。在 PEPS 层面上，可以通过将基态张量中的两个替换为支集包含在标记该激发的相应幂等元的像中的张量，来创建一对任意子激发，详细内容请参阅文献~\cite{Bultinck2017}。

这里考虑的（无扭）量子双重模型可以在许多方面进行推广。一方面，可以考虑该模型的\emph{有扭}（twisted）版本，其中 MPO 对称性构成 $G$ 的一个反常表示，其异常取值于第三上同调群 $H^3(G,\rU(1))$，正如在\cref{sec:SPT1d}中所描述的那样。对此，我们推荐参阅文献~\cite{Buerschaper2013,Bultinck2016,Williamson2017}。

在\cref{sec:SN}中，我们将讨论 Levin 和 Wen 在文献~\cite{Levin2004}中首次构造的所谓弦网模型（string-net models）的 PEPS 表示。这是一类以球面融合范畴为输入并实现所有非手征拓扑序 $\mc Z(\mc D)$ 的模型。在这些模型中获得由 $\mc Z(\mc D)$ 所描述的环面基态和任意子激发的一种方法，恰好是通过与 $\mc D$ 相关联的管代数，这与上述构造完全一致。对于输入 $\mc D=\Vect_G^\omega$，相应的 Levin-Wen 模型与有扭量子双重模型相重合。
\subsection{第四讲小结} 
第四讲介绍了三个新主题：构成特殊类量子自旋链可积性基础的 MPO 形式的对称性；MPS 向 PEPS 的高维推广，以及参数化有能隙高维晶格模型低能态的相应流形；以及满足所谓\emph{穿透对称性}（pulling through symmetry）的 MPO 对称性，该对称性表征了此类系统的不同有能隙（拓扑）相。关于可积性，我们论证了 Yang-Baxter 方程是 MPS 基本定理的一个推论。如果这种观点能够引向更一般的可积模型类别（例如满足非平庸范畴对称性的模型，参见第五讲），那将是令人振奋的。

我们展示了二维晶格系统中的 SPT 相由其结合子编码了 3-上循环的非平庸 MPO 对称性所表征，从而导出了二维 SPT 相的分类。然而，在不涉及全局对称性的情况下，PEPS 的穿透对称性为在二维系统中实现拓扑物态开辟了全新的途径。我们以量子双重模型的形式展示了内禀拓扑序晶格系统的存在，并证明了这些系统相对于纠缠自由度表现出对称性破缺。这是第一个启示：如果用底层 PEPS 表示的对称性来描述多体物态，Landau 范式便可以得到拯救——不同的相在根本不同的方式下实现张量的对称性。